% has to be loaded outside of a frame to work!
\maketitle

\begin{frame}[plain, noframenumbering]{Licença}
    \begin{columns}
        \begin{column}{0.5\textwidth}
            O texto e as figuras desses slides possuem uma
            \href{https://creativecommons.org/licenses/by-nc-sa/4.0/deed.pt}{Licença
            Creative Commons
            Atribuição-NãoComercial-CompartilhaIgual 4.0 Internacional (CC BY-NC-SA 4.0)}
            \vfill
            \centering
            \vspace{1em}
            \includegraphics[width=0.2\columnwidth]{CC_NC_SA.png}
        \end{column}
        \begin{column}{0.5\textwidth}
            \centering
            \includegraphics[width=0.8\columnwidth]{qrcode.png}
        \end{column}
    \end{columns}
    \vfill
\end{frame}

\begin{frame}{Truque da Corrente Linear\footnote{\textit{Linear Chain Trick}.} (TCL)}
	O Truque da Corrente Linear (TCL) \parencite{hurtadoGeneralizationsLinearChain2019, hurtadoProcedureDerivingNew2020, andoHowFastLinear2020}
	é usado em modelos epidemiológicos compartimentais para
	modelar tempos de espera em transições de compartimentos usando
	subcompartimentos com uma distribuição \textcolor{blue}{Erlang}\footnote{mais informações
	sobre a distribuição Erlang na \href{https://en.wikipedia.org/wiki/Erlang_distribution}{Wikipedia}.}.
	\vfill
	A distribuição Erlang é uma soma de $k$ distribuições exponenciais
	independentes.
	\vfill
	Sem subcompartimentos, a distribuição do tempo de espera se torna uma
	distribuição \textcolor{blue}{Exponencial}.
\end{frame}

\begin{frame}{Distribuição Exponencial vs Erlang}
    \begin{columns}
        \begin{column}{0.5\textwidth}
	    \begin{figure}
	    \centering
	    \begin{tikzpicture}[scale=0.5]
       	    	\begin{axis}[every axis plot, line width=2pt,
            	ylabel=FDP,
            	xlabel={$X$},
            	domain=0:5,samples=200,
            	axis x line*=bottom, % no box around the plot, only x and y axis
            	axis y line*=left, % the * suppresses the arrow tips
            	enlarge x limits=true, % extend the axes a bit
            	]

		\addplot [blue] {exponential(1)};
            	\addlegendentry{$\lambda=1$}
	    	\end{axis}
            \end{tikzpicture}
	    \caption{Distribuição Exponencial}
   	    \end{figure}
        \end{column}
        \begin{column}{0.5\textwidth}
	    \begin{figure}
            \centering
	    \begin{tikzpicture}[scale=0.5]
       	    	\begin{axis}[every axis plot, line width=2pt,
            	ylabel=FDP,
            	xlabel={$X$},
            	domain=0:5,samples=200,
            	axis x line*=bottom, % no box around the plot, only x and y axis
            	axis y line*=left, % the * suppresses the arrow tips
            	enlarge x limits=true, % extend the axes a bit
            	]

		\addplot [red] {erlang(1,1)};
            	\addlegendentry{$\lambda=1, k=1$}
		\addplot [yellow] {erlang(1,2)};
            	\addlegendentry{$\lambda=1, k=2$}
		\addplot [green] {erlang(1,3)};
            	\addlegendentry{$\lambda=1, k=3$}
	    	\end{axis}
            \end{tikzpicture}
	    \caption{Distribuição Erlang}
   	    \end{figure}
        \end{column}
    \end{columns}
    \vfill
\end{frame}

\begin{frame}[allowframebreaks]{Referências}
	\printbibliography
\end{frame}
